\documentclass[12pt,a4paper]{article}

\usepackage{fontawesome5}
\providecommand{\faPencilAlt}{\faIcon{pencil-alt}}

\usepackage{styles/couleurs}
\usepackage{styles/mise_en_page}
\usepackage{styles/boites}
\usepackage{styles/maths}
\usepackage{styles/figures}
\usepackage{styles/evaluations}

\usepackage{enumitem}
\usepackage{array}
\usepackage{tabularx}
\usepackage{amsmath}
\usepackage{tikz}

\usetikzlibrary{calc,angles,quotes}

\setlist[itemize]{leftmargin=1.4em,itemsep=0.15em,topsep=0.2em}
\setlist[enumerate]{leftmargin=*,itemsep=0.25em,topsep=0.25em}

\begin{document}

\enteteinterro[Calculatrice autorisée]{4e}{Calculer une longueur avec le cosinus}{18 min}{12}

\exo{4}[Calculer]
Dans le triangle \(DEF\) rectangle en \(E\),

\[
DF=9\text{ cm}
\qquad
\widehat D=34^\circ.
\]

Calcule \(DE\), arrondi au dixième de centimètre.
Rédige entièrement. \pts{4}

\lignes{7}

\exo{4}[Calculer]
Dans le triangle \(GHI\) rectangle en \(H\),

\[
GH=6{,}5\text{ m}
\qquad
\widehat G=49^\circ.
\]

Calcule l'hypoténuse \(GI\), arrondie au dixième de mètre.
Rédige entièrement. \pts{4}

\lignes{7}

\exo{2}[Raisonner]
Sans refaire le calcul, explique pourquoi le résultat suivant est impossible :

\[
\text{côté adjacent}=12{,}4\text{ cm}
\]

alors que l'hypoténuse mesure \(10\) cm. \pts{2}

\lignes{3}

\exo{2}[Communiquer]
Complète :

\[
\text{Pour calculer le côté adjacent, je }\ldots
\]

\[
\text{Pour calculer l'hypoténuse, je }\ldots
\]

\lignes{2}

\end{document}
